% Sections 1--3, drafting pass 16 October 2026.
% Empirical scope follows the merged Sections 4--6 review snapshot.
\section{Introduction}
\label{sec:introduction}

An agent buying on behalf of a firm must decide more than whether it can
place an order. It must choose whom to trust, whether another test is worth
its cost, and when to commit before an offer expires. A tenant faces the
same problem when inspecting a home; a negotiator faces it when exchanging
price for contractual risk. In these settings, completing a transaction
can destroy value, while refusing it can be the right decision. We use
\emph{economic-agent readiness} to mean an agent's ability to obtain value
for a principal within a stated mandate, given limited information and
other parties' interests.

Evaluating this ability requires a reference for what the agent could
have achieved. A good realized outcome can follow a poor decision, and a
shortfall from a fully informed optimum can reflect information the agent
never had. The counterpart also matters: a strategy that succeeds against
a cooperative seller may fail against a different policy. Existing work
provides interactive task evaluation and economic decision benchmarks
\citep{merrill2026terminalbench,fish2025econevals,zhang2026termsbench}.
The problem we address is how to make information, counterpart behavior
and reference values explicit across different economic cases, then test
whether the resulting scores measure the intended decisions.

AERead provides a framework for constructing and evaluating such cases.
Each case specifies the parties, their observations and legal actions,
the world distribution, and checks on behavior and outcomes. A shared
runtime executes play and records evidence. Case-specific references give
scores an economic interpretation: a feasible policy, an optimal policy
under the deciding seat's information, or a labelled full-information
upper bound. Results retain their native units and reference types.

CASEFIGUREPLACEHOLDER

AERead combines seven properties. Prior benchmarks have some of them,
TERMS-Bench five in full; AERead has six in full and the seventh, a
shortcut-screening admission rule, in three of its cases (Table~2.1). The first three
concern a single world, the next two a case's population of worlds, and the
last two the runtime (Figure~1.1). Each is built and was applied to the
cases reported here.
\begin{enumerate}
\item \textbf{Value against a reference with stated information.} Scores
are economic values in each case's own units, measured against a feasible
policy, the best policy given the deciding seat's information, or a
labelled full-information bound. Each reference is played through the
environment on every admitted world (Section~\ref{sec:qc}).
\item \textbf{The counterpart is part of the case.} Other parties are seats
with their own information and interests, played by scripted policies,
another model or the same model, and each claim names its counterpart. In
the Data-center case, one agent's earnings lead over another is 42\% of the
value available in a deal head to head and 7\% against a fixed scripted
counterpart
(Section~\ref{sec:cases}).
\item \textbf{Typed shortfall.} The gap to the reference is split into
observed action classes, such as tests not taken, worse terms, harmful
agreements and negotiation fees, so a result states where value was lost
as well as how much (Section~\ref{sec:cases}).
\item \textbf{World generators with admission.} A case is a generator of
worlds, and its packs regenerate from their seeds. An admission rule,
computed before any agent plays, keeps only worlds in which simple
shortcuts lose a stated margin: the Data-center comparison used 80 worlds
admitted from 5,436 scanned. Housing's v15 pack admits a world only if the best informed tenant
beats a simple inspect-then-sign rule by \$25, and the reported Procurement pack only if a rule
that plans one newcomer's test exactly pays the premium while the best policy beats paying it by
10\%: 20 worlds from 5,164 scanned (Section~\ref{sec:cases}).
\item \textbf{Quality control over populations of worlds.} Cases are
qualified through typed gates; simple-rule ceilings are scored beside the
agents and briefs are audited. Of 25 defects selected from our record, 14
were invisible in any single world (Section~\ref{sec:qc}).
\item \textbf{Scores recomputed from sealed evidence.} The kernel hashes a
plan before execution, seals the events, and replays each receipt to
recompute its scores without calling a model; no model judges a score
(Section~\ref{sec:runtime}).
\item \textbf{One runtime for several activities.} Families supply only the
economic rules. Four maintainer-built cases and a reimplementation of three
TERMS-Bench counterpart families run on the same kernel
(Section 5.3).
\end{enumerate}
A researcher studying another economic decision therefore writes only its
rules, worlds and reference: execution, sealing, replay, counterpart seats
and the quality-control gates come with the runtime, so a new case starts
with the same evidence and checks as the existing ones.
Sections 3 and 4 describe the case design and the runtime, Section 5 the
cases, and Section 6 the quality control. The cases are maintainer-built
worked examples; Housing passes four of the five qualification gates and
none passes all five; Section 7
states what their results do and do not establish.

\section{Related work}
\label{sec:related-work}

Table~2.1 compares the closest benchmarks on the seven properties of
Section~\ref{sec:introduction}; the paragraphs below discuss them by theme.

COMPARISONTABLEPLACEHOLDER

\paragraph{Economically valuable work.}
GDPval selects professional tasks by occupation and evaluates deliverables
through expert comparison \citep{patwardhan2025gdpval}. Agents' Last Exam
(ALE) evaluates expert-authored professional workflows with task-specific
verification \citep{sun2026agentslastexam}. AERead evaluates decisions
under a principal's mandate, with responsive counterparts and payoff
references that specify the deciding party's information.

\paragraph{Interactive agent evaluation.}
Terminal-Bench verifies final environment state and supplies human-written
solutions \citep{merrill2026terminalbench}. AppWorld combines an application
engine with state-based evaluation, including collateral-change checks
\citep{trivedi2024appworld}. τ²-bench lets agent and user act on shared
state \citep{barres2025tau2}.

\paragraph{Economic decisions and negotiation.}
EconEvals covers procurement, scheduling, pricing and economic choice
\citep{fish2025econevals}. NegotiationArena studies bargaining tactics
\citep{bianchi2024negotiationarena}; SidConArena combines negotiation,
production and auctions \citep{feng2026sidconarena}. AgenticPay evaluates
buyer--seller bargaining under private reservation values, scoring
agreement, surplus allocation and speed \citep{liu2026agenticpay}.
TERMS-Bench specifies counterpart types, policies and payoffs to diagnose
negotiation through oracle-reference gaps \citep{zhang2026termsbench}; our
case specifications follow its decomposition into environment, protocol
and counterpart policy (Appendix~A), and Section 5.3 reimplements three of
its counterpart families on AERead's runtime.
AERead retains case-specific payoff units and distinguishes full-information
bounds from information-constrained policy references, with QC of world
distributions and measurement paths.

\paragraph{Shared environments and evaluation quality.}
OpenSpiel and PettingZoo provide common interfaces for games and
multi-agent environments \citep{lanctot2019openspiel,terry2021pettingzoo}.
AERead similarly separates shared execution from case semantics
(Appendix H). Terminal-Bench combines executable reference solutions,
automated checks, human review and adversarial exploit testing
\citep{merrill2026terminalbench}. The Agentic Benchmark Checklist
identifies task and reward defects that distort performance
\citep{zhu2025abc}. Our QC analysis tests whether prompts, controls and
counterpart assignments support valid comparisons across distributions
of economic worlds.

\section{Task design and construction}
\label{sec:task-design}

Three of AERead's properties are fixed when a case is designed: the
reference and the information it may use, the counterpart seats, and the
world generator with its admission rule.

\subsection{Cases, worlds and seats}
\label{sec:task-formulation}

A \emph{case} defines an economic decision problem; a \emph{world} is
one instance of that design, with particular valuations, private types,
costs and deadlines. A \emph{seat} is a party's position in the
interaction, with its own objective, information and permitted actions.
A policy occupies a seat and chooses actions from its observations and
available history. The family implementation supplies the environment,
observation rules, legal actions, transitions and verifier.

For example, Housing places a tenant among four listings with hidden
quality. The tenant can inspect, make offers and decide whether to sign;
landlord replies and outside demand affect the opportunities that remain.
Inspection changes both information and cost. The verifier therefore
needs to distinguish a legal action, a favorable realized outcome and a
good decision given the tenant's information.

In a controlled comparison, the evaluated policy occupies one seat and
counterparts follow fixed policies. Cross-play and self-play change those
assignments and answer different questions. Each comparison names the
seat, counterpart configuration and world distribution to which it
applies. Appendices A--D give the case specifications, counterpart
policies, references and measures.

\subsection{Selecting economic decisions}

We organize coverage around activities in which a party's choices affect
the value of a transaction \citep{li2026activities}. A candidate case
needs a consequential decision, a counterpart with its own interests and
information, and an outcome that can be evaluated inside the world.
The intended coverage spans business purchases, labour contracts, capital
and corporate transactions, household purchases, and public purchases.
Contractual terms, risk allocation and disputes enter as decisions within
these transactions.

The four examples cover business and household purchases. They were
selected to exercise different references and counterpart configurations;
they are not a representative sample of economic activity. Economic
stakes motivate which decisions to study, but the reported scores are
case-specific quantities, without a cross-family aggregate or an estimate
of economy-wide gains.

\subsection{Constructing and qualifying a case}
\label{sec:case-construction}

A contribution begins with a decision and an interpretation of success.
The author specifies the principal's objective and constraints, what each
seat can observe, which actions commit resources, and how counterpart
policies respond. The brief must expose the rules needed to make the
decision while preserving the intended private information. These
choices define the mechanism that a score is meant to test.

The author then provides a world generator or fixed pack, an admission
rule, and a reference whose information access is explicit. Admission
screens each candidate world before any agent plays and before hidden
types are drawn. Computing exactly from the public world and the declared
prior, it keeps a world only if simple shortcuts, such as accepting the
default terms or ignoring private information, fall short of the reference
by a stated margin. Whether the reference can be reached is a separate
check: executable controls play it through the actual environment and
verifier on every admitted world, alongside simple policies that reveal
shortcuts or ceiling effects that admission missed. The pack must also be
valid, free of duplicates, and split into disjoint development and holdout
ranges. A full-information bound is labelled separately from a policy
the evaluated seat could execute. The case declares score units,
direction and treatment of invalid actions and incomplete episodes.

Before a confirmatory comparison, the campaign freezes the worlds,
seat assignments, prompts, agent settings, budgets and analysis. Its
evidence package should connect each reported measure to the retained
actions and outcomes, the scoring implementation, and a command that
reproduces it. Section 4 describes the shared execution path; Appendix H
records where the current drivers and headline analyses extend beyond
the receipt's coverage.

This is the contribution and qualification standard. Implementation
support and completed checks vary across the four examples, as Section 6
and Appendix F document. Cases can be useful for development before
qualification, provided their reports identify missing checks, preserve
failed attempts, and distinguish analyses declared before execution from
later corrections. The current examples demonstrate the workflow and
its failure modes; they do not establish successful adoption by external
contributors.
